\chapter{Additional Results}
\label{app:additional-results}

This appendix provides supplementary ablation results that support Chapter~\ref{chap:results}. All reported values are tied to retained public artefacts. The complete OOD detection results (including FPR@95\%TPR) are presented in the main text (Table~\ref{tab:results:main-auroc}).

\section{Separation Loss Ablation: Detailed Results}
\label{app:additional-results:separation}

Table~\ref{tab:appD:lambda-ablation} shows the effect of the separation loss coefficient $\lambda$ on within-CIFAR OOD detection performance. All results use $K{=}50$ MC trials and difference scoring. The $\lambda \in \{0.00, 0.01, 0.02\}$ entries are mean $\pm$ std over three seeds.

\begin{table}[htbp]
  \centering
  \caption{Separation loss ablation: AUROC (\%) on within-CIFAR split as a function of $\lambda$. All conditions use $K{=}50$ and difference scoring. The $\lambda \in \{0.00, 0.01, 0.02\}$ rows report three seeds (42/123/456).}
  \label{tab:appD:lambda-ablation}
  \begin{tabular}{lcc}
    \toprule
    $\lambda$ & AUROC (\%) $\uparrow$ & $\Delta$ vs.\ $\lambda{=}0$ \\
    \midrule
    0.00 (no separation loss) & $92.52 \pm 11.07$ & --- \\
    0.01 & $98.82 \pm 0.06$ & $+6.30$ \\
    0.02 & $99.03 \pm 0.07$ & $+6.51$ \\
    \bottomrule
  \end{tabular}
\end{table}

The results show that the main effect of a non-zero $\lambda$ is reduced seed sensitivity. The peak is achieved at $\lambda{=}0.02$, which is the best/primary reported setting. The $\lambda{=}0.01$ checkpoint is retained for the $K$, timestep-sampling, and scoring-strategy ablations because those runs were completed before the final three-seed sweep and were not rerun at $\lambda{=}0.02$. Larger $\lambda$ values risk over-constraining the model, which we leave to future work.


\section{Monte Carlo Trials Ablation: Detailed Results}
\label{app:additional-results:mc-trials}

Table~\ref{tab:appD:k-ablation} characterises the accuracy--efficiency frontier as a function of the number of MC timestep trials $K$.

\begin{table}[htbp]
  \centering
  \caption{$K$ ablation: AUROC (\%) and relative inference cost on within-CIFAR split using the seed-42 $\lambda{=}0.01$ checkpoint. Reference cost is the discriminative baseline (ResNet-18 forward pass) measured on an NVIDIA GeForce RTX 2080 Ti. The measured wall-clock timings are 972.9\,s for $K{=}10$ and 4861.1\,s for $K{=}50$ per 10K images.}
  \label{tab:appD:k-ablation}
  \begin{tabular}{lclp{4cm}}
    \toprule
    $K$ & AUROC (\%) $\uparrow$ & Inference Cost & Notes \\
    \midrule
    1 & 91.0 & not re-timed in final consistency pass & Low latency \\
    10 & 98.2 & $68.6\times$ discriminative & Recommended trade-off \\
    50 & 98.64 & $492.4\times$ discriminative & Default (within-CIFAR) \\
    \bottomrule
  \end{tabular}
\end{table}

The performance gain from $K{=}1$ to $K{=}10$ is substantial ($+7.2$ pp), while the gain from $K{=}10$ to $K{=}50$ is marginal ($+0.44$ pp). This suggests that $K{=}10$ is the practical sweet spot for most deployment scenarios. For the external OOD evaluation, $K{=}100$ was used to maximise statistical stability of the reported means; $K{=}50$ produces similar AUROC values on those datasets.


\section{Scoring Strategy: Difference vs.\ ID-Only}
\label{app:additional-results:scoring}

Table~\ref{tab:appD:scoring} compares contrastive difference scoring ($S = e_0(x) - e_1(x)$, where $c = 0$ is ID and $c = 1$ is non-ID) against single-sided ID-only scoring ($S = e_0(x)$) on selected OOD datasets.

\begin{table}[htbp]
  \centering
  \caption{Scoring strategy comparison using the seed-42 $\lambda{=}0.01$ checkpoint at $K{=}50$: AUROC (\%) for contrastive difference, ratio, and ID-only scoring.}
  \label{tab:appD:scoring}
  \begin{tabular}{lcccc}
    \toprule
    Dataset & Difference & Ratio & ID-Only & Best -- ID-Only \\
    \midrule
    Within-CIFAR & 98.69 & 98.62 & 78.30 & $+20.39$ \\
    SVHN & 94.13 & 96.06 & 20.2 & $+75.86$ \\
    \bottomrule
    \multicolumn{5}{@{}p{0.95\linewidth}@{}}{\footnotesize All values come from the retained seed-42 $\lambda{=}0.01$ scoring ablation JSON at $K{=}50$. The main text reports difference scoring because it is the default reproducible scorer.}
  \end{tabular}
\end{table}

The catastrophic failure of ID-only scoring on SVHN (20.2\% AUROC, below random chance on a balanced evaluation) occurs because SVHN images have lower absolute reconstruction error under the ID model than many within-CIFAR OOD images, making the absolute magnitude uninformative. Contrastive scoring resolves this by measuring the relative advantage of the ID model, providing a signal that is robust to absolute reconstruction magnitude differences. Ratio scoring is competitive here, but the thesis retains difference scoring as the primary reproducible scorer for consistency across tables and raw-score evaluation.
